\cleardoublepage
\phantomsection
\pdfbookmark{Sommario}{Sommario}
\begingroup
\let\clearpage\relax
\let\cleardoublepage\relax
\let\cleardoublepage\relax

\chapter*{Sommario}

Il presente documento descrive il lavoro svolto durante il periodo di stage, della durata di circa trecento ore, presso l'azienda Sia Informatica s.r.l.
\\
Il progetto nasce da due esigenze concrete legate a \emph{Sia.Core}, la piattaforma
condivisa da tutti i prodotti aziendali, distribuita su decine di installazioni
cliente: da un lato, la necessità di una rete di sicurezza che permetta di evolvere
il \emph{framework} senza rompere ciò che già funziona; dall'altro, la necessità di
rendere visibile il comportamento dell'applicazione in produzione, dove un
malfunzionamento può manifestarsi lontano dalla sua causa reale e con un ritardo
difficile da imputare.\\
La risposta a queste esigenze si articola in due filoni paralleli. Il primo riguarda
la \textbf{telemetria e l'osservabilità}: è stata progettata un'architettura completa
basata su \emph{OpenTelemetry} per la raccolta distribuita di metriche, \emph{trace}
e \emph{log}, con \emph{Grafana} come piattaforma unificata di visualizzazione e
analisi. Il secondo riguarda il \textbf{testing automatico}: è stata introdotta
un'infrastruttura di test articolata su \emph{unit test}, \emph{integration test}
e test \emph{end-to-end}, volta a garantire la correttezza del
\emph{framework} al variare del codice.\\
Un elemento trasversale a entrambi i filoni è l'utilizzo strutturato dell'intelligenza
artificiale come assistente allo sviluppo: agenti AI specializzati per modulo, guidati da
\emph{skill} procedurali e connessi direttamente agli strumenti di sviluppo tramite
\gls{mcp}, hanno reso la propagazione di telemetria
e test ai nuovi moduli un processo replicabile e omogeneo, riducendo le allucinazioni
e aumentando la qualità del codice generato.

\endgroup

\vfill
